\elementopretextual{Resumo}{resumo}
\begingroup
\singlespacing
\noindent
O envelhecimento cutâneo envolve alterações da matriz extracelular e da pigmentação que motivam a investigação de interações entre compostos naturais e proteínas relacionadas a esses processos. Este trabalho comparou, por \textit{docking} molecular, monossacarídeos, derivados de açúcares e fenólicos selecionados para prospecção no contexto das SPE de \textit{Tetraselmis chuii}, utilizando ligantes comerciais e de referência como comparadores. Foram avaliados 34 ligantes, divididos em dois conjuntos de 17, em seis estruturas humanas: três metaloproteinases de matriz (MMPs), MMP-9 (1GKC), MMP-1 (1HFC) e MMP-2 (8H78); elastase (1H1B); hialuronidase 1 (2PE4); e proteína relacionada à tirosinase 1 (TYRP1, 5M8P). A seleção bibliográfica distinguiu evidências em biomassa, parede celular, exopolissacarídeos e vesículas. O AutoDock Vina foi empregado com pontuação Vina para elastase e hialuronidase e com a função AutoDock4 (AD4) e a parametrização AutoDock4Zn (AD4Zn) para os outros quatro receptores. A análise reuniu 204 pontuações FEB (do inglês \textit{Free Energy of Binding}, energia livre de ligação) por combinação receptor--ligante, complementadas por visualização no PyMOL, contatos identificados pela ferramenta \textit{Protein--Ligand Interaction Profiler} (PLIP) e comparação conformacional com referências cristalográficas disponíveis. Entre os compostos selecionados para prospecção, destacaram-se glicose-6-sulfato em MMP-9 e MMP-1, galactose-6-sulfato em MMP-2, Dha (ácido \textit{3-deoxy-D-lyxo-hept-2-ulopyranosaric acid}) em elastase e hialuronidase e ácido ferúlico em TYRP1. Os líderes comerciais apresentaram pontuações mais negativas nos seis receptores. Os resultados organizam os perfis de pontuação e contatos e orientam a priorização de compostos para investigação dirigida. A caracterização química da SPE, seguida de ensaios bioquímicos e celulares, permitirá avaliar a presença desses compostos e seu potencial de aplicação cutânea.

\vspace{\baselineskip}
\noindent\textbf{Palavras-chave:} \textit{docking} molecular; \textit{Tetraselmis chuii}; substâncias poliméricas extracelulares; envelhecimento cutâneo; triagem computacional.
\endgroup
